\section{Introduction}
\label{sec:introduction}

Let $r\ge3$ and $n\ge r$ be integers, write $[n]=\{1,\ldots,n\}$,
and put $E_{n,r}=\binom{[n]}r$.
An element of $\{0,1\}^{E_{n,r}}$ represents a simple $r$-uniform
hypergraph. A \emph{hypergraph property} is a Boolean function on this
space that is invariant under every permutation of $[n]$.
For $0<p<1$, give the coordinates the independent Bernoulli$(p)$ law,
and denote expectation and variance by $\E_p$ and $\Var_p$.

\noindent\textit{Remark (notation).} $p$ denotes the common Bernoulli bias, with $p\in[0,1]$. In Theorem~\ref{thm:main}, Lemma~\ref{lem:low}, Lemma~\ref{lem:block}: $0<p<1$. In Corollary~\ref{cor:threshold}: In $q(p)=\E_p f$ and the definition of $p_a$, $p$ ranges over $[0,1]$, including the endpoint product laws. $A$ is a subset of the current coordinate set indexing a Fourier basis function or coefficient. In Lemma~\ref{lem:low}: In the difference identity $\widehat h(A\cup\{i\})=\sigma\widehat{\Delta_i h}(A)$, $A\subseteq J\setminus\{i\}$. In Lemma~\ref{lem:block}: In the expansion of $\widehat{f_y}(A)$, $A\subseteq F_B$.


For an edge $e$, let $x_{e\leftarrow b}$ be $x$ with its $e$-coordinate
set to $b$. Define the \emph{pivotal influence} of $e$ and the
\emph{total influence} of $f$ by
\[
 I_{e,p}(f)=\Pr_p\bigl(f(x_{e\leftarrow1})\ne
                         f(x_{e\leftarrow0})\bigr),
 \qquad I_p(f)=\sum_{e\in E_{n,r}}I_{e,p}(f).
\]
The probability in the first expression is over the other coordinates;
it does not include an additional factor of $p(1-p)$.
All logarithms are natural, and $n$ always counts vertices.

\noindent\textit{Remark (notation).} $x$ denotes a configuration in $\{0,1\}^{J}$ for the coordinate set on which the function being evaluated is defined. In Theorem~\ref{thm:main}, Corollary~\ref{cor:threshold}: For evaluations of $f$, the coordinate set is $E_{n,r}$. In Lemma~\ref{lem:block}: The coordinate set is $E=E_{n,r}$ for $f$, and $F_B$ for the restricted function $f_y$. $b$ is locally bound, with its domain determined by the expression in which it occurs. In Theorem~\ref{thm:main}, Lemma~\ref{lem:block}, Corollary~\ref{cor:threshold}: In the coordinate replacement notation $x_{e\leftarrow b}$, $b\in\{0,1\}$ is the replacement value. In Lemma~\ref{lem:low}: In the one-coordinate norm calculation, $b\in\mathbb R$ is an arbitrary scalar. In the product induction, $g=a+b\chi_i$ uses a real-valued coefficient function $b$ on $\{0,1\}^{J\setminus\{i\}}$.


\begin{theorem}\label{thm:main}
For every integer $r\ge3$ there is a finite constant $C_r$ such that,
for every integer $n\ge r$, every $0<p<1$, and every hypergraph property
$f:\{0,1\}^{E_{n,r}}\to\{0,1\}$,
\begin{equation}\label{eq:main}
 \Var_p(f)\le \frac{C_r}{(\log n)^{r/(r-1)}}I_p(f).
\end{equation}
No monotonicity assumption is required.
\end{theorem}

The theorem measures how vertex symmetry forces sensitivity to edge
changes. If the probability of a property stays bounded away from
zero and one, its total influence must grow at least as
$(\log n)^{r/(r-1)}$.

A property is \emph{increasing} if changing coordinates from $0$ to
$1$ cannot decrease its value, and \emph{nontrivial} if it is not
constant. For a nontrivial increasing property, put $q(p)=\E_p f$
and define $p_a=\inf\{p\in[0,1]:q(p)\ge a\}$ for $0<a<1$.
For fixed $0<\varepsilon<1/2$, its \emph{threshold width} is
$p_{1-\varepsilon}-p_\varepsilon$.
Friedgut and Kalai conjectured that this width is at most
$O_{r,\varepsilon}((\log n)^{-r/(r-1)})$ for every increasing
$r$-uniform hypergraph property~\cite[Section~5, p.~2999]{FK96}.
Their formulation counts edge coordinates rather than vertices;
for fixed $r$, $\log\binom nr\sim r\log n$ as $n\to\infty$.
Corollary~\ref{cor:threshold} proves this conjectured bound by
integrating~\eqref{eq:main}, with an explicit dependence on
$\varepsilon$.

\noindent\textit{Remark (notation).} The letter $a$ is locally bound, with its domain and meaning determined by the calculation in which it occurs. In Corollary~\ref{cor:threshold}: $a\in(0,1)$ is the probability level indexing $p_a$. In Lemma~\ref{lem:low}: In the one-coordinate calculation, $a\in\mathbb R$ is a coefficient; in the product induction $g=a+b\chi_i$, it is a real-valued function on $\{0,1\}^{J\setminus\{i\}}$.


\paragraph{Context and attribution.}
Kahn, Kalai, and Linial~\cite{KKL88} established a fundamental link
between Fourier analysis and the influences of Boolean functions.
For a function invariant under a transitive group on $M$ coordinates,
their theorem gives a total-influence lower bound of order
$\Var_{1/2}(f)\log M$ at the uniform measure.
Bourgain, Kahn, Kalai, Katznelson, and Linial~\cite{BKKKL92}
extended the influence method to product spaces.
Friedgut and Kalai~\cite{FK96} used influence inequalities to prove
sharp thresholds for increasing families with transitive symmetry.
For hypergraph properties, the induced action on the
$\binom nr$ edge coordinates is transitive. The richer structure of
vertex permutations, however, contains more information than this
single transitivity statement.

Bourgain and Kalai~\cite[Corollary~4.1]{BK97} studied how group
actions on sets of coordinates improve influence estimates.
For fixed $r$ and every $\eta>0$, their result gives a lower bound
of order $\Var_{1/2}(f)(\log n)^{r/(r-1)-\eta}$ for the total
influence of a hypergraph property at the uniform measure.
Their influence results already allow nonmonotone functions.
Kelman, Kindler, Lifshitz, Minzer, and Safra obtained further
estimates for group-invariant functions~\cite{KKLMS20}.
For graph properties, these give
$I_{1/2}(f)\ge c\Var_{1/2}(f)(\log n)^2/(\log\log n)^2$
for sufficiently large $n$; see the arXiv version cited
in~\cite[Section~7.4]{KKLMS20}.
The present proof uses vertex symmetry through restrictions: all
edges touching a vertex block remain free, and only edges wholly
outside the block are fixed. This is the method of the companion
manuscript \emph{A Sharp Threshold Bound for Monotone Graph
Properties}~\cite[Theorem~1.2 and Sections~2--3]{Graph2026}, which establishes
the corresponding variance bound with denominator $(\log n)^2$
for graph properties, uniformly in $p$ and without monotonicity.
We reproduce its Fourier estimate and restriction argument in full.
The additional combinatorial step replaces the graph degree estimate
by its $r$-uniform counterpart, producing the exponent $r/(r-1)$.

\paragraph{Proof strategy.}
Put $\sigma=\sqrt{p(1-p)}$ and $\alpha=r/(r-1)$.
The Fourier expansion for the Bernoulli product measure is indexed by
sets $S$ of hyperedges. Its squared nonconstant coefficients sum to
the variance; weighting them by $|S|$ gives $\sigma^2I_p(f)$.
Thus the influence controls high degrees. The task is to bound the
remaining mass using symmetry.

\noindent\textit{Remark (notation).} $S$ denotes a subset of the current coordinate set, used as a Fourier index or, for hypergraph coordinates, as an edge support. In Lemma~\ref{lem:low}: $S\subseteq J$, and the low-degree sum ranges over $1\le|S|\le k$. In Theorem~\ref{thm:main}, Lemma~\ref{lem:block}: $S\subseteq E_{n,r}$ indexes the Fourier expansion of $f$. In Lemma~\ref{lem:capture}: $S\subseteq E_{n,r}$ is a fixed support with $1\le|S|\le L$, considered before $B$ is sampled.


For large $n$, choose a uniformly random block $B$ of $m$ vertices,
with $m$ about $\sqrt n$, and leave free all edges meeting $B$.
Every permutation within $B$ survives an assignment to the other
coordinates, and every free-edge orbit has at least $m$ elements.
Section~\ref{sec:fourier} bounds the restricted Fourier mass up to
a degree $k$ proportional to $\sigma\log m$ by $m^{-1/4}$ times
the restricted total influence. Its linear dependence on influence
allows averaging over the assigned coordinates.

\noindent\textit{Remark (notation).} $m$ is a fixed real parameter with $m\ge1$. In Lemma~\ref{lem:low}: $m\ge2$ is real and is the parameter in the individual-influence bound $I_i(h)\le I(h)/m$ and the cutoff $k=\sigma\log m/64$. In Lemma~\ref{lem:block}: $m$ is an integer with $r+1\le m\le n$, and $|B|=m$. In Lemma~\ref{lem:capture}: $m$ is an integer with $1\le m\le n$, and $|B|=m$. In Theorem~\ref{thm:main}: In the proof for $n\ge N_r$, $m=\lfloor\sqrt n\rfloor$ is the block size. $k$ is a real degree cutoff; comparisons of integer cardinalities with $k$ require no rounding. In Lemma~\ref{lem:low}: $k=\sigma\log m/64$ with real $m\ge2$; the low-degree sum is empty if $k<1$. The auxiliary projection inequalities are established for every real cutoff $k\ge0$. In Lemma~\ref{lem:block}: $k=\sigma\log m/64$ with integer $r+1\le m\le n$. In Lemma~\ref{lem:capture}: $k>0$ is arbitrary, and $L=(k/r)^{r/(r-1)}$. In Theorem~\ref{thm:main}: In the branch $n\ge N_r$, $m=\lfloor\sqrt n\rfloor$ and $k=\sigma\log m/64>0$.


Regard a nonempty Fourier index $S$ as an $r$-uniform hypergraph
with $s$ edges. It has a vertex of positive degree at most
$r s^{(r-1)/r}$. For $s\le L=(k/r)^\alpha$, this degree is at
most $k$. The event that $B$ meets the union of these edges only
at that vertex has probability comparable to $m/n$ and leaves
between one and $k$ support edges free. After averaging over the
assignments and $B$, the restricted total influence is at most
$(rm/n)I_p(f)$, so the two factors of $m/n$ cancel.
Section~\ref{sec:restriction} thereby controls the original
Fourier mass up to degree $L$.

\noindent\textit{Remark (notation).} $L=(k/r)^\alpha$ is a positive real support-size cutoff. In Lemma~\ref{lem:capture}: $k>0$ is an arbitrary cutoff, so $L$ is not tied to $\sigma\log m/64$. In Theorem~\ref{thm:main}: In the branch $n\ge N_r$, $m=\lfloor\sqrt n\rfloor$, $k=\sigma\log m/64$, and hence $L=(\sigma\log m/(64r))^\alpha$. $s=|S|$ is the number of edges in the support under consideration. In Lemma~\ref{lem:capture}: $1\le s\le L$.


Since $L$ is of order $(\sigma\log n)^\alpha$, the remaining
mass is at most $\sigma^2I_p(f)/L$, with coefficient of order
$\sigma^{2-\alpha}(\log n)^{-\alpha}$. The positive exponent
$2-\alpha$ makes this bound uniform in $p$.
Section~\ref{sec:completion} completes the estimate and derives
the threshold consequence.
